\section{Results for Rectina m1 against Reader v3}\label{sec:results-MB}

Rectina m1 is the checkpoint at the end of the first stage. Table~\ref{tab:comparison-MB} gives, in the format of the table of the previous section, the three measures of Reader v3 and Rectina m1 on the fourteen units.

\MBHonestSentence{}

\begin{table*}[!t]
\caption{Three measures of \MBModelName{} against Reader v3 on the fourteen units, target a: the AUC, the AUC on thin
strokes and false ink at 80 per cent recall. The cell of \MBModelName{} gives in parentheses the signed difference
against Reader v3. The rows Mean and Median are taken \MBStatScope; the
block of the community labels is information only.}\label{tab:comparison-MB}
\centering
\definecolor{rgreen}{RGB}{46,139,74}\definecolor{rred}{RGB}{192,57,43}
{\small\setlength{\tabcolsep}{2.2pt}
\noindent\begin{tabular}{>{\raggedright\arraybackslash}p{3.0cm}>{\raggedleft\arraybackslash}p{1.5cm}>{\raggedleft\arraybackslash}p{3.0cm}>{\raggedleft\arraybackslash}p{1.5cm}>{\raggedleft\arraybackslash}p{3.0cm}>{\raggedleft\arraybackslash}p{1.5cm}>{\raggedleft\arraybackslash}p{3.0cm}}
\hline
unit & \multicolumn{2}{c}{AUC} & \multicolumn{2}{c}{\shortstack{AUC on\\thin strokes}} & \multicolumn{2}{c}{\shortstack{false ink at 80\\per cent recall}} \\
 & {\footnotesize Reader v3} & {\footnotesize Rectina m1} & {\footnotesize Reader v3} & {\footnotesize Rectina m1} & {\footnotesize Reader v3} & {\footnotesize Rectina m1} \\
\hline
\multicolumn{7}{@{}p{\dimexpr 16.5cm+30.8pt\relax}@{}}{\rule{0pt}{2.6ex}\textit{Control units outside the training segments, organisers' labels: 0343 (approved letters, 116 keV); 0841 (113 keV)}} \\
0343 & 0.9078 & \textcolor{rgreen}{$\blacktriangle$}\,\textbf{0.9105} {\scriptsize (+0.0026)} & 0.9246 & 0.9246 {\scriptsize (+0.0000)} & \textbf{0.0690} & \textcolor{rred}{$\blacktriangledown$}\,0.0705 {\scriptsize (+0.0014)} \\
0841 ag1 & \textbf{0.8513} & \textcolor{rred}{$\blacktriangledown$}\,0.8395 {\scriptsize (-0.0119)} & \textbf{0.8648} & \textcolor{rred}{$\blacktriangledown$}\,0.8519 {\scriptsize (-0.0129)} & \textbf{0.2153} & \textcolor{rred}{$\blacktriangledown$}\,0.2530 {\scriptsize (+0.0376)} \\
0841 ag2 & 0.8297 & \textcolor{rgreen}{$\blacktriangle$}\,\textbf{0.8449} {\scriptsize (+0.0153)} & 0.8474 & \textcolor{rgreen}{$\blacktriangle$}\,\textbf{0.8600} {\scriptsize (+0.0126)} & 0.2681 & \textcolor{rgreen}{$\blacktriangle$}\,\textbf{0.2631} {\scriptsize (-0.0050)} \\
0841 w00 & \textbf{0.9055} & \textcolor{rred}{$\blacktriangledown$}\,0.8826 {\scriptsize (-0.0229)} & \textbf{0.9221} & \textcolor{rred}{$\blacktriangledown$}\,0.8997 {\scriptsize (-0.0223)} & \textbf{0.1444} & \textcolor{rred}{$\blacktriangledown$}\,0.1939 {\scriptsize (+0.0495)} \\
\multicolumn{7}{@{}p{\dimexpr 16.5cm+30.8pt\relax}@{}}{\rule{0pt}{2.6ex}\textit{Control unit outside the training segments, community labels, information only: 1447 (116 keV)}} \\
1447 w058 & \textbf{0.8610} & \textcolor{rred}{$\blacktriangledown$}\,0.8491 {\scriptsize (-0.0118)} & \textbf{0.8965} & \textcolor{rred}{$\blacktriangledown$}\,0.8865 {\scriptsize (-0.0100)} & \textbf{0.2665} & \textcolor{rred}{$\blacktriangledown$}\,0.2813 {\scriptsize (+0.0148)} \\
1447 w060 & 0.7842 & \textcolor{rgreen}{$\blacktriangle$}\,\textbf{0.7958} {\scriptsize (+0.0116)} & 0.8276 & \textcolor{rgreen}{$\blacktriangle$}\,\textbf{0.8429} {\scriptsize (+0.0153)} & 0.4095 & \textcolor{rgreen}{$\blacktriangle$}\,\textbf{0.3770} {\scriptsize (-0.0325)} \\
\multicolumn{7}{@{}p{\dimexpr 16.5cm+30.8pt\relax}@{}}{\rule{0pt}{2.6ex}\textit{Held-out segments of training scrolls, our validation labels: 0009B and 0343P (116 keV), 0500P2 (113 keV)}} \\
0009B$\cdot$172639 & 0.8456 & \textcolor{rgreen}{$\blacktriangle$}\,\textbf{0.9259} {\scriptsize (+0.0803)} & 0.8515 & \textcolor{rgreen}{$\blacktriangle$}\,\textbf{0.9303} {\scriptsize (+0.0788)} & 0.3114 & \textcolor{rgreen}{$\blacktriangle$}\,\textbf{0.0960} {\scriptsize (-0.2154)} \\
0009B$\cdot$125754 & 0.9562 & \textcolor{rgreen}{$\blacktriangle$}\,\textbf{0.9600} {\scriptsize (+0.0038)} & 0.9636 & \textcolor{rgreen}{$\blacktriangle$}\,\textbf{0.9663} {\scriptsize (+0.0027)} & 0.0460 & \textcolor{rgreen}{$\blacktriangle$}\,\textbf{0.0455} {\scriptsize (-0.0006)} \\
0009B$\cdot$133208 & 0.8551 & \textcolor{rgreen}{$\blacktriangle$}\,\textbf{0.8773} {\scriptsize (+0.0222)} & 0.8672 & \textcolor{rgreen}{$\blacktriangle$}\,\textbf{0.8873} {\scriptsize (+0.0201)} & 0.2774 & \textcolor{rgreen}{$\blacktriangle$}\,\textbf{0.2203} {\scriptsize (-0.0571)} \\
0009B$\cdot$140325 & \textbf{0.9040} & \textcolor{rred}{$\blacktriangledown$}\,0.8784 {\scriptsize (-0.0256)} & \textbf{0.9110} & \textcolor{rred}{$\blacktriangledown$}\,0.8860 {\scriptsize (-0.0250)} & \textbf{0.1574} & \textcolor{rred}{$\blacktriangledown$}\,0.1908 {\scriptsize (+0.0334)} \\
0500P2$\cdot$055229 & \textbf{0.8971} & \textcolor{rred}{$\blacktriangledown$}\,0.8605 {\scriptsize (-0.0366)} & \textbf{0.9074} & \textcolor{rred}{$\blacktriangledown$}\,0.8698 {\scriptsize (-0.0376)} & \textbf{0.1360} & \textcolor{rred}{$\blacktriangledown$}\,0.2854 {\scriptsize (+0.1494)} \\
0500P2$\cdot$055230 & 0.7881 & \textcolor{rgreen}{$\blacktriangle$}\,\textbf{0.7984} {\scriptsize (+0.0103)} & 0.7980 & \textcolor{rgreen}{$\blacktriangle$}\,\textbf{0.8078} {\scriptsize (+0.0098)} & 0.3998 & \textcolor{rgreen}{$\blacktriangle$}\,\textbf{0.3541} {\scriptsize (-0.0458)} \\
0343P$\cdot$233748-0 & 0.8340 & \textcolor{rgreen}{$\blacktriangle$}\,\textbf{0.8584} {\scriptsize (+0.0244)} & 0.8398 & \textcolor{rgreen}{$\blacktriangle$}\,\textbf{0.8629} {\scriptsize (+0.0231)} & 0.3059 & \textcolor{rgreen}{$\blacktriangle$}\,\textbf{0.2669} {\scriptsize (-0.0390)} \\
0343P$\cdot$172054-1 & 0.7573 & \textcolor{rgreen}{$\blacktriangle$}\,\textbf{0.7981} {\scriptsize (+0.0407)} & 0.7669 & \textcolor{rgreen}{$\blacktriangle$}\,\textbf{0.8035} {\scriptsize (+0.0366)} & 0.4336 & \textcolor{rgreen}{$\blacktriangle$}\,\textbf{0.3679} {\scriptsize (-0.0657)} \\
\hline
Mean & 0.8555 & \textcolor{rgreen}{$\blacktriangle$}\,\textbf{0.8628} {\scriptsize (+0.0073)} & 0.8706 & \textcolor{rgreen}{$\blacktriangle$}\,\textbf{0.8771} {\scriptsize (+0.0065)} & 0.2458 & \textcolor{rgreen}{$\blacktriangle$}\,\textbf{0.2333} {\scriptsize (-0.0125)} \\
\hline
Median & 0.8532 & \textcolor{rgreen}{$\blacktriangle$}\,\textbf{0.8595} {\scriptsize (+0.0062)} & 0.8660 & \textcolor{rgreen}{$\blacktriangle$}\,\textbf{0.8779} {\scriptsize (+0.0119)} & 0.2673 & \textcolor{rgreen}{$\blacktriangle$}\,\textbf{0.2580} {\scriptsize (-0.0093)} \\
\hline
\end{tabular}}

\par\smallskip
{\footnotesize\raggedrightAUC: the rank AUC of ink against blank pixels. \\
AUC on thin strokes: the rank AUC of the central 80 per cent of each stroke against blank pixels at least 207 um from any label, the band between not scored. \\
false ink at 80 per cent recall: the share of blank pixels marked as ink at the threshold that recalls 80 per cent of the ink pixels, where less is better. \\
For the controls the AUC is taken against the far negatives of the bench; for the held-out units against all blank pixels of the validation mask. \\
The better value of each pair in bold; a green triangle pointing up before \MBModelName{} where it is better, a red triangle pointing down where it is worse. \\
A pair whose two values are equal at four decimals has no bold value and no triangle.
\par}
\end{table*}

On the eight held-out units the mean score is \MBHeldMeanReader{} for Reader v3 and \MBHeldMeanModel{} for \MBModelName{}; \MBModelName{} scores above Reader v3 on \MBHeldAbove{} of the 8
units, the median difference is \MBHeldMedianDiff{} and the mean difference is \MBHeldMeanDiff{}\MBLowHeld. On the six control units \MBModelName{} scores above Reader v3 on \MBStatsCtrlKText{} with a median difference of \MBStatsCtrlMedianDiff{}. On the four controls with organisers' labels (ag1, ag2, w00, 0343), a subset of the six, the split following from who made the labels of each unit, a fact about the data and not a choice made on the scores (both counts are given because the split was not fixed before the scores), the mean score is \MBCtrlMeanReader{} for Reader v3 and \MBCtrlMeanModel{} for
\MBModelName{}; \MBModelName{} scores above Reader v3 on
\MBCtrlAbove{} of the \MBCtrlN{} units, the median difference is \MBCtrlMedianDiff{} and the mean difference is
\MBCtrlMeanDiff{}\MBLowCtrl.

\MBFalseInkParagraph{}

\MBStatsParagraph{} \MBSeedParagraph{}
